\section{Introduction}\label{sec:introduction}

Let $m$ denote Lebesgue measure on $\R$. A set $A\subseteq\R$ is
\emph{measure universal} if every Lebesgue-measurable set of positive
measure contains a nontrivial affine copy
\[
 x+sA=\{x+sa:a\in A\},\qquad x\in\R,\quad s\in\R\setminus\{0\}.
\]
The Erd\H{o}s similarity conjecture \cite[Problem~4.33.7*]{Erdos1974Problems}
asserts that no infinite set is measure universal. We prove its
geometric-progression case. For
$q\in(0,1)$, write
\[
 G_q=\{q^n:n\in\N,\ n\ge1\}.
\]

\begin{theorem}\label{thm:main}
For every $q\in(0,1)$ and $\eta\in(0,1)$, there is a compact set
$E_{q,\eta}\subseteq[0,1]$ with $m(E_{q,\eta})>1-\eta$ such that
\[
 \text{for every }x\in\R\text{ and }s\in\R\setminus\{0\},
 \quad x+sq^n\notin E_{q,\eta}\text{ for some integer }n\ge1.
\]
\end{theorem}

The avoiding set may depend on $q$. In particular, taking $q=1/2$
gives a compact set of arbitrarily large measure in $[0,1]$ avoiding
every signed affine copy of the dyadic sequence. The theorem treats
all geometric ratios, while the conjecture for arbitrary infinite
sets is a separate question.

\subsection{Background and related work}

Every finite pattern is measure universal. Indeed, if $F$ is finite
and $A$ is measurable with finite positive measure, continuity of
translations in $L^1$ gives
\[
 m\Bigl(\bigcap_{a\in F}(A-sa)\Bigr)\longrightarrow m(A)>0
 \qquad(s\to0).
\]
For an arbitrary positive-measure set, apply this observation to a
bounded subset of positive measure. The conjecture asks whether this
property always fails for infinite patterns.

Falconer~\cite{Falconer1984} and Eigen~\cite{Eigen1985} proved
nonuniversality for positive decreasing sequences $a_n\to0$ with
$a_{n+1}/a_n\to1$. Humke and Laczkovich
\cite[Theorems~1--2 and Corollary~1]{HumkeLaczkovich1998} developed a
finite covering characterization and further criteria in terms of
relative gaps. Kolountzakis~\cite[Theorem~3]{Kolountzakis1997} gave
a probabilistic criterion using large finite subsets with sufficiently
large normalized minimum gap. In the translation-invariant form of
Chleb\'ik~\cite[Theorem~15]{Chlebik2015}, a bounded infinite set is
nonuniversal if it contains finite subsets $a_1>\cdots>a_m$, of
arbitrarily large cardinality, for which
\[
 -\log\left(\frac{\min_{i<m}(a_i-a_{i+1})}{a_1-a_m}\right)=o(m).
\]
Geometric progressions do not satisfy this criterion. For distinct
powers of $q$, we have $a_{m-1}\le q^{m-2}a_1$ and
$a_1-a_m\ge(1-q)a_1$. Their normalized minimum gap is therefore at
most $q^{m-2}/(1-q)$, whose negative logarithm cannot be sublinear
in $m$.

Additive structure supplies another family of results.
Bourgain~\cite{Bourgain1987} proved that the sum of three infinite
sets is nonuniversal. Kolountzakis~\cite{Kolountzakis1997} treated
certain double sums, including $G_{1/2}+G_{1/2}$.
Mora Cuellar, Iosevich, Kulkarni, Rojas Aravena and
Yavicoli~\cite[Theorem~2]{MoraCuellarEtAl2026Twofold} proved that the
sum or difference of a geometric null sequence and any infinite set
is nonuniversal. Avoiding such a larger pattern does not imply
avoidance of the single geometric progression in Theorem~\ref{thm:main}.

Other results concern different classes of patterns or notions of
largeness. Iosevich, Kulkarni, Mora Cu\'ellar, Rojas Aravena and
Yavicoli~\cite[Theorem~1.5]{IKMRY2026Rajchman} prove nonuniversality
for sets supporting a probability measure whose Fourier transform
tends to zero at infinity. Such a measure is atomless, so the
hypothesis excludes countable sets
\cite[Proposition~4.1]{IKMRY2026Rajchman}.
Cruz, Lai and Pramanik~\cite[preprint, Corollary~1.2 and
Proposition~4.4]{CruzLaiPramanik2023} construct sets of Hausdorff
dimension one avoiding certain null sequences together with their
limit point; the subsets establishing that dimension have Lebesgue
measure zero. In a different direction, the theorem of Feng, Lai and
Xiong~\cite[preprint, Theorem~1.1]{FengLaiXiong2024} implies that
every $G_q$ embeds into every measurable set of positive measure by
a bi-Lipschitz map. The restriction to affine maps is therefore
substantive.

The proof below belongs to the probabilistic approach of
Kolountzakis~\cite{Kolountzakis1997}. Random cells, fixed-center
scale discretization, and integration of exceptional-center
probabilities also appear in Chleb\'ik~\cite[Section~5]{Chlebik2015}
and Kolountzakis--Papageorgiou~\cite[Section~3.1]{KolountzakisPapageorgiou2025};
the latter paper treats unbounded sequences. Solymosi~\cite[Section~3]{Solymosi2011Sparse}
and Tom~\cite{Tom2015} report incomplete dyadic constructions based on
random cells; Tom also considers nested dyadic intervals and tests
independent of the dilation. Our contribution is the
finite routing construction and its local control of the scale
count. We explain that mechanism next and give a self-contained proof.


\subsection{The proof mechanism}

We construct the open complement of the avoiding set. The intermediate
goal is an open, $1$-periodic set of arbitrarily small density that meets
$x+tG_q$ for every $x\in\R$ and every $t\in[1,2]$. A summable union
of its dilations and reflections then handles every nonzero scale;
Section~\ref{sec:periodic} gives this reduction. The remaining proof
uses only finite probability spaces and elementary measure theory.

For a fixed center and a fixed scale, many independent random cell
tests make a missed copy unlikely. The difficulty is to control all
scales at once. Recording every possible change at the finest grid
of a large construction can cost more than the failure probability
can afford. We arrange the tests on a finite ordered tree so that each
group of tests uses only the grids on one edge and in its child subtree.
Other parts of the tree may use much finer grids without increasing
the scale count for that group.

Each point is routed to a leaf by independent binary tables evaluated
at its grid cells. A node chooses the first child whose table returns
one, and chooses its last child by default if all earlier tables return
zero. The leaf has a separate table whose entries equal one with a
small probability $p$; this final table determines membership in the
random set and keeps its expected density equal to $p$.

We assign a consecutive interval of sequence indices to each edge,
in the order that traverses an edge before its entire child subtree.
Gaps between these intervals make earlier routing decisions unchanged
by the translations associated with a later interval, except at a
small set of centers. Within the current interval, the translated
points use distinct entries of its own table. At the first default
on the center's route, the nondefault children therefore supply
many independent opportunities for a hit. Taking the tree sufficiently
deep makes the probability of having no default small.

The decisive length condition makes an edge's index interval at least
as long as the following gap and its child subtree together. Every
grid needed by a test from that interval then occurs within twice its
length. Thus the number of scale representatives depends on the
length of one interval, rather than on the resolution of the entire
tree. The branching number can be chosen so that the independent-test
failure probability beats this scale count.

For general $q$, we retain nested dyadic grids but choose their
resolutions comparable to $q^{-b}/(1-q)$ at index $b$. This gives
both the nesting needed to preserve earlier decisions and the
separation needed for fresh random entries. The branching number is
allowed to depend on $q$. These two adjustments let the same finite
routing mechanism accommodate any fixed geometric ratio.

Finally, the construction may miss copies at a small closed set of
centers. An open neighborhood of those centers repairs every such
copy, because $x+tq^n$ tends to $x$. This converts the estimate for
exceptional centers into a statement valid at every center.
Section~\ref{sec:windows} constructs the grids and index intervals;
Section~\ref{sec:routing} proves the conditional independence of the
tests; and Section~\ref{sec:scales} controls all normalized scales
and completes the open-neighborhood repair.
